%!TEX program = xelatex
\PassOptionsToPackage{table,svgnames}{xcolor}
\PassOptionsToPackage{french}{translator}
%\documentclass[aspectratio=169,11pt]{beamer}
\documentclass[11pt]{beamer}

%----------------------------------------
% Packages
%----------------------------------------
%\usepackage[utf8]{inputenc}
\usepackage[T1]{fontenc}
\usepackage{translator}
\usepackage{lmodern}
\usepackage{hyperref}
\usepackage{xcolor}
\usepackage{listings}
\usepackage{amsmath}
\usepackage{amssymb}
\usepackage{mathrsfs}
\usepackage{array}
\usepackage{tabularx}
\usepackage{multirow}
\usepackage[justification=centering]{caption}
\usepackage{float}
\usepackage{standalone}
\usepackage{import}
% PGF-TikZ
\usepackage{pgf}
\usepackage{tikz}
\usepackage{pgf-umlsd}
\usepackage{pgfgantt}
% Personnal packages
\usepackage{packages/MagicListings}
\usepackage{packages/StandardLibraryCDefinition}

%----------------------------------------
% Theme
%----------------------------------------

\usetheme[illustration=cover2]{utbm}

%----------------------------------------
% Informations
%----------------------------------------

\title{Développement et fiabilisation d'un module de désassemblage interne}
%\subtitle{Direction Générale de l'Armement}
\author{Maxime Pinard}
\institute[UTBM]{Université de Technologie de Belfort Montbéliard}
\date{1er mars 2018}

%\keywords{}
\subject{Développement et fiabilisation d'un module de désassemblage interne}
%\logo{\includegraphics[width=0.12\textwidth]{logos/DGA}}

%----------------------------------------
% Configurations
%----------------------------------------

\graphicspath{{figures/}}

\AtBeginSection
{
	\begin{frame}[plain]
		\utbmtitle{\insertsectionhead}
	\end{frame}
}

%----------------------------------------
% Listings configuration
%----------------------------------------

\lstalias[gendbg]{C}[GenDbg]{C}
\lstdefinelanguage[gendbg]{C}{
	language={[StandardLibrary]C},
	morekeywords=[1]{
		__cdecl
	},
	morekeywords=[2]{
		DWORD
	},
	morekeywords=[2]{
		AsmBankInfo_T,
		AsmAddressSpaceInfo_T,
		AsmAddressTypeInfo_T,
		AsmGroupRegisterInfo_T,
		AsmDataInfo_T,
		AsmDecodedInstruction_T,
		AsmInstructionInfo_T,
		AsmModuleInfo_T,
		CPUCtx_T,
		GenDbgHelperAsmInfo_T,
		MemoryAddress_T,
		MemoryArea_T,
		ViewCtx_T,
		AsmModuleFnIdx_T,
		MIPS_RegisterId_T
	},
	morekeywords=[2]{
		cs_insn,
		cs_detail,
		cs_x86,
		cs_arm64,
		cs_arm,
		cs_m68k,
		cs_mips,
		cs_ppc,
		cs_sparc,
		cs_sysz,
		cs_xcore,
		cs_tms320c64x,
		cs_mips_op,
		mips_reg,
		mips_op_mem
	},
	morekeywords=[3]{
		CS_MNEMONIC_SIZE
	}
}


%----------------------------------------
% Figures configuration
%----------------------------------------

\usetikzlibrary{shapes}
\usetikzlibrary{arrows.meta}
\usetikzlibrary{calc}

\definecolor{bg_color}{RGB}{250,250,229}

\colorlet{color1}{cyan!50}
\colorlet{color2}{red!30!green!40}
\colorlet{color3}{orange!50}
\colorlet{color4}{violet!60!blue!55}

\newganttlinktype{bartobardown}{
	\ganttsetstartanchor{south east}
	\ganttsetendanchor{north west}
	\draw [/pgfgantt/link] (\xLeft, \yUpper) -- (\xRight, \yLower);
}
\newganttlinktype{bartobarup}{
	\ganttsetstartanchor{north east}
	\ganttsetendanchor{south west}
	\draw [/pgfgantt/link] (\xLeft, \yUpper) -- (\xRight, \yLower);
}
\newganttlinktype{milestonetobardown}{
	\ganttsetstartanchor{south}
	\ganttsetendanchor{north west}
	\draw [/pgfgantt/link] (\xLeft, \yUpper) -- (\xRight, \yLower);
}
\newganttlinktype{bartomilestonedown}{
	\ganttsetstartanchor{south east}
	\ganttsetendanchor{north}
	\draw [/pgfgantt/link] (\xLeft, \yUpper) -- (\xRight, \yLower);
}


%----------------------------------------
% Document
%----------------------------------------
\begin{document}
	\begin{frame}[plain,noframenumbering]
		\titlepage
	\end{frame}
	\begin{frame}{Sommaire}
		\tableofcontents
	\end{frame}
	\section{Présentation de l'entreprise}
		\begin{frame}
			\frametitle{Du Ministère des Armées à la DGA MI}
			\centering\resizebox{\textwidth}{!}{\import{figures/}{uncover_slide_Ministere_Defense_a_DGA_MI.tex}}
			%----------
			% MinArm: 265 000 personnes
			% DGA: 9 700 personnes
			%----------
			% Missions DGA:
			% - Équiper les forces armées
			% - Préparer l’avenir
			% - Promouvoir les exportations d’armement
		\end{frame}
		\begin{frame}
			\frametitle{La DGA MI}
			\centering\resizebox{\textwidth}{!}{\import{figures/}{Hierarchie_DGA_MI_a_VIM.tex}}
			% Ancien CELAR
			% Conception, expertise, essais, évaluation: systemes d'informations, informatique, communications, spacial
		\end{frame}
		\begin{frame}
			\frametitle{VIM/VSE}
			\begin{block}{Analyse dynamique}
				\begin{itemize}
					\item Gendbg
						\begin{itemize}
							\invisible{\item Développement et fiabilisation du module de désassemblage MIPS}
						\end{itemize}
				\end{itemize}
			\end{block}
			\pause
			\begin{block}{Analyse statique}
				\begin{itemize}
					\item IDA
						\begin{itemize}
							\item YaCo
								\begin{itemize}
									\invisible{\item Portage en C++ et amélioration de la gestion des évènements IDA et du dépôt Git}
								\end{itemize}
						\end{itemize}
				\end{itemize}
			\end{block}
		\end{frame}
	\section{Sujets du stage}
		\begin{frame}
			\frametitle{Sujets du stage}
			\begin{block}{Analyse dynamique}
				\begin{itemize}
					\item Gendbg
						\begin{itemize}
							\item[\alert{$\blacktriangleright$}]<2-> Développement et fiabilisation du module de désassemblage MIPS
						\end{itemize}
				\end{itemize}
			\end{block}
			\begin{block}{Analyse statique}
				\begin{itemize}
					\item IDA
						\begin{itemize}
							\item YaCo
								\begin{itemize}
									\item[\alert{$\blacktriangleright$}]<3-> Portage en C++ et amélioration de la gestion des évènements IDA et du dépôt Git
								\end{itemize}
						\end{itemize}
				\end{itemize}
			\end{block}
		\end{frame}
	\section{Travail réalisé}
		\begin{frame}
			\frametitle{Utilisation}
			\centering\resizebox{\textwidth}{!}{\import{figures/}{GenDbg_architecture_minimale.tex}}
		\end{frame}
		\begin{frame}
			\frametitle{GUI}
			\centering\resizebox{\textwidth}{!}{\import{figures/}{uncover_GenDbg_GUI_Qt_Annotations.tex}}
		\end{frame}
		\begin{frame}
			\frametitle{GUI: vue code}
			\centering\resizebox{\textwidth}{!}{\import{figures/}{uncover_GenDbg_GUI_Qt_Vue_code_Annotations.tex}}
		\end{frame}
		\begin{frame}
			\frametitle{Le framework modulaire}
			\centering\resizebox{\textwidth}{!}{\import{figures/}{uncover_GenDbg_modules.tex}}
		\end{frame}
		\begin{frame}
			\frametitle{But du stage}
			\centering\resizebox{\textwidth}{!}{\import{figures/}{uncover_GenDbg_modules_impl.tex}}
		\end{frame}
		\begin{frame}
			\frametitle{Assemblage désassemblage}
			\centering\resizebox{\textwidth}{!}{\import{figures/}{Capstone_Keystone_utilisation_basique.tex}}
		\end{frame}
	\section{Résultats}
		\begin{frame}
			\frametitle{Résultats}
			\begin{block}{Ancien module MIPS}
				\begin{itemize}
					\item 8 200 lignes
				\end{itemize}
			\end{block}
			\begin{block}{Nouveau module MIPS}
				\begin{itemize}
					\item 1 600 lignes pour le module
					\item 1 500 lignes pour les tests
				\end{itemize}
			\end{block}
			\uncover<2->{
				\begin{block}{Dépendances}
					\begin{itemize}
						\item \ 56 600 lignes pour Capstone
						\item 100 000 lignes pour Keystone
					\end{itemize}
				\end{block}
			}
		\end{frame}
	\section{Conclusion}
		\begin{frame}
			\frametitle{Conclusion}
			\begin{block}{Personnellement}
				\begin{itemize}
					\item Découverte de nouveaux métiers
					\item Expérience des contraintes de travail dans un contexte de très haute sécurité
				\end{itemize}
			\end{block}
			\begin{block}{Professionnellement}
				\begin{itemize}
					\item Travail sur des gros projets à durée de vie longue
					\item Acquisition de nouvelles compétences
				\end{itemize}
			\end{block}
		\end{frame}
		\begin{frame}[plain]
			\utbmclosingframe{Avez-vous des questions?}
		\end{frame}
		\subsection*{Stage}
			\begin{frame}\label{planning}
				\frametitle{Planning}
				\hspace*{-10pt}%
				\centering\resizebox{0.9\paperwidth}{!}{\import{figures/}{Planning.tex}}
			\end{frame}
\end{document}
